% Figure: the plant recirculating-power fraction against fusion gain Q. A
% power plant must run its own magnets, heating, and balance of plant from a
% slice of its own output; that recirculating fraction falls as 1/Q (times a
% conversion and heating-efficiency factor), so a plant needs Q well above the
% break-even Q=1 to sell net electricity. The curve marks the ITER Q=10 point
% and the plant-relevant Q~30-40 band.
\begin{figure}[htbp]
\centering
\begin{tikzpicture}
\begin{axis}[
  width=0.82\textwidth, height=0.55\textwidth,
  xlabel={fusion gain $Q=P_{\mathrm{fus}}/P_{\mathrm{heat}}$},
  ylabel={recirculating power fraction $f_{\mathrm{rec}}$},
  xmin=1, xmax=50, ymin=0, ymax=1.0,
  grid=both, grid style={enggray!20},
  every axis plot/.append style={thick},
]
  % f_rec = 1/(eta_th * (1 + Q/5) * eta_heat) approx; use eta_th=0.4,
  % eta_heat=0.7, and gain in electricity M ~ (1 + Q/5) for alpha+neutron.
  % f_rec = 1 / (0.4 * 0.7 * (1 + Q/5))
  \addplot[engblue, domain=1:50, samples=80]
    {1/(0.4*0.7*(1 + x/5))};
  % ITER Q=10 marker
  \addplot[only marks, mark=*, engred, mark size=2pt] coordinates {(10,1.19)};
  \draw[dashed, enggray] (axis cs:1,1) -- (axis cs:50,1);
  \node[enggray, font=\scriptsize, anchor=south east] at (axis cs:50,1.0)
    {break-even in electricity};
  \node[engred, font=\scriptsize, anchor=west] at (axis cs:10,0.85)
    {ITER $Q=10$};
  % plant band
  \draw[engorange, thick, |-|] (axis cs:30,0.15) -- (axis cs:40,0.15);
  \node[engorange, font=\scriptsize, anchor=north] at (axis cs:35,0.15)
    {plant band};
\end{axis}
\end{tikzpicture}
\caption[Recirculating power against fusion gain]{The fraction of its own
output a fusion plant must recirculate to run its magnets, heating, and
balance of plant, drawn against the fusion gain $Q$. The estimate takes a
thermal-conversion efficiency near $0.4$, a heating-injection efficiency near
$0.7$, and an electricity multiplication set by the alpha and neutron energy,
so the recirculating fraction falls roughly as the inverse of the gain. At the
$Q=10$ ITER point the plant would still spend more than its whole output on
itself, which is why ITER is a physics-gain demonstration and not a power
plant; a plant that sells electricity needs $Q$ in the thirties or forties,
where the recirculating fraction drops to a fifth or less of the output. The
gap between the $Q=10$ that demonstrates the burn and the $Q\sim 40$ that
sells power is the engineering distance the fusion programme still has to
cover.}
\label{fig:vol04ch13:fusion-gain}
\end{figure}
